\chapter{Tag 3 — Hamming-Codes}

Stub-Kapitel für M1. Inhalte folgen in M3.

\clearpage

\begin{bleistiftuebung}{1}{Aufgabe direkt am Seitenanfang}
Dies ist \emph{Szenario 5} aus dem Briefing: die erste Aufgabe steht
ganz oben auf einer Seite. Das vertikale Alignment der Marginalie darf
durch das fehlende vorhergehende Element nicht gestört werden.

Datenwort: \texttt{1011}. Bestimme das Codewort des
\((7,4)\)-Hamming-Codes.
\end{bleistiftuebung}

\begin{werkzeugcheck}{Buntstifte für die drei Kreise}
Halte drei Buntstifte in Orange, Blau und Grün bereit — du wirst sie
im nächsten Block für das Hamming-Venn-Diagramm brauchen.
\end{werkzeugcheck}

Eine kleine TikZ-Skizze, damit später bestätigt werden kann, dass die
Diagramme weiterhin sauber rendern:

\begin{center}
\begin{tikzpicture}[scale=0.7]
  \def\R{1.6}
  \draw[thick, draw=cBleistift, fill=cBleistift, fill opacity=0.25]
        (-0.8, 0.5) circle (\R);
  \draw[thick, draw=cPython,    fill=cPython,    fill opacity=0.25]
        ( 0.8, 0.5) circle (\R);
  \draw[thick, draw=cWerkzeug,  fill=cWerkzeug,  fill opacity=0.25]
        ( 0.0,-0.9) circle (\R);
\end{tikzpicture}
\end{center}
